\documentclass[11pt, reqno]{amsart}

\usepackage[spanish]{babel}
\usepackage[LGR, T1]{fontenc}
\usepackage[utf8]{inputenc}

../LaTeX/general.tex
% \input{../graphics.tex}

\makeatletter
\def\emailaddrname{\textit{Correo electrónico}}
\def\subtitle#1{\gdef\@subtitle{#1}}
\def\@subtitle{}

% Metadata
\def\logo#1{\gdef\@logo{#1}}
\def\@logo{}
\def\institution#1{\gdef\@institution{#1}}
\def\@institution{}
\def\department#1{\gdef\@department{#1}}
\def\@department{}
\def\professor#1{\gdef\@professor{#1}}
\def\@professor{}
\def\course#1{\gdef\@course{#1}}
\def\@course{}
\def\coursecode#1{\gdef\@coursecode{#1}}
\def\@coursecode{}

\renewcommand{\maketitle}{
\begin{center}
	\small
	\renewcommand{\arraystretch}{1.2}
	\begin{tabular}{cp{.37\textwidth}p{0.44\textwidth}}
		% \hline
		\multirow{5}{*}{\includegraphics[height=2.0cm]{\@logo}}
	  & \multicolumn{2}{c}{ \makecell{{\bfseries \@institution} \\ \@department} } \\
	  % & \multicolumn{2}{c|}{{\bfseries\@institution} \\ \@department} \\
	  \cline{2-3}
	  & \textbf{Profesor:} \@professor & \textbf{Ayudante:} \authors \\
	  % \cline{2-3}
	  & \textbf{Curso:} \@course & \textbf{Sigla:} \@coursecode \\
	  % \cline{2-3}
	  & \multicolumn{2}{l}{ \textbf{Fecha:} \@date } \\
	  % \hline
	\end{tabular}
	\\[\baselineskip]
	% {}
	% \vspace{2\baselineskip}
	{\bfseries\Large\@title}
	\ifx\@subtitle\@empty\else
		\\[1ex]
		\large\mdseries\@subtitle
	\fi
\end{center}
}
\makeatother

\usepackage{multirow, makecell}

\usepackage[
	reversemp,
	letterpaper,
	% marginpar=2cm,
	% marginsep=1pt,
	margin=2.3cm
]{geometry}
\usepackage{fontawesome}
% \makeatletter
% \@reversemargintrue
% \makeatother

% Símbolos al margen, necesitan doble compilación
\newcommand{\hard}{\marginnote{\faFire}}
\newcommand{\hhard}{\marginnote{\faFire\faFire}}

% Dependencias para los teoremas
\usepackage{xifthen}
\def\@thmdep{}
\newcommand{\thmdep}[1]{
	\ifthenelse{\isempty{#1}}
	{\def\@thmdep{}}
	{\def\@thmdep{ (#1)}}
}
\newcommand{\thmstyle}{\color{thm}\sffamily\bfseries}

% ===== Estilos de Teoremas ==========
\newtheoremstyle{axiomstyle}
	{0.3cm}
	{0.3cm}
	{\normalfont}
	{0.5cm}
	{\bfseries\scshape}
	{:}
	{4pt}
	{\thmname{#1}\thmnote{ #3}\thmnumber{ (#2)}}
\newtheoremstyle{styleC}
	{0.5cm}
	{0.5cm}
	{\normalfont}
	{0.5cm}
	{\bfseries}
	{:}
	{4pt}
	{\thmname{#1\textrm{\@thmdep}}\thmnumber{ #2}\thmnote{ (#3)}}

% ====== Teoremas (sin borde) ===========
\theoremstyle{axiomstyle}
\newtheorem*{axiom}{Axioma}

% ====== Teoremas (sin borde) ==================
\theoremstyle{styleC}
\newtheorem{thm}{Teorema}[section]
\newtheorem{mydef}[thm]{Definición}
\newtheorem{prop}[thm]{Proposición}
\newtheorem{cor}[thm]{Corolario}
\newtheorem{lem}[thm]{Lema}
\newtheorem{con}[thm]{Conjetura}
\newtheorem*{sol}{Solución}
\newtheorem*{obs}{Observación}
\newtheorem*{ex}{Ejemplo}

\newtcbox{bluebox}[1][]{enhanced jigsaw, 
  sharp corners,
  frame hidden,
  nobeforeafter,
  listing only,
  #1} % comando para crear cajas de colores

\expandafter\let\expandafter\oldproof\csname\string\proof\endcsname
\let\oldendproof\endproof
\renewenvironment{proof}[1][\proofname]{%
  \oldproof[\scshape Demostración:]%
}{\oldendproof} % comando para redefinir la caja de la demostración
\newenvironment{hint}[1][\proofname]{%
  \oldproof[\scshape Pista:]%
}{\oldendproof} % comando para redefinir la caja de la demostración

% colores utilizados
\definecolor{numchap}{RGB}{249,133,29}
\definecolor{chap}{RGB}{6,129,204}
\definecolor{sec}{RGB}{204,0,0}
\definecolor{thm}{RGB}{106,176,240}
\definecolor{thmB}{RGB}{32,31,31}
\definecolor{part}{RGB}{212,66,66}

% ====== Diseño de los titulares ===============
\usepackage[explicit]{titlesec} % para personalizar el documento, la opción <<explicit>> hace que el texto de los titulares sea un objeto interactuable

\titleformat{\subsection}[runin]
	{\bfseries}
	{\textrm{\S}\thesubsection}
	{1ex}
	{#1.}

\setlist[enumerate,1]{label=\arabic*., ref=\arabic*} % Enumerate standards

\usepackage{tikz}
\usetikzlibrary{babel,cd}
% \DeclareMathOperator\Gr{Gr}

\DeclareMathOperator\pr{pr}

\title{Conexión del tipo Galois}
\date{\DTMdate{2026-10-22}}

\author{José Cuevas Barrientos}
\email{josecuevasbtos@uc.cl}
\urladdr{https://josecuevas.xyz/teach/2026-2-ayud/}

\logo{../puc_negro.png}
\institution{Pontificia Universidad Católica de Chile}
\department{Facultad de Matemáticas}
\course{Topología Algebraica}
\coursecode{MAT2850}
\professor{Mauricio Bustamante}

\begin{document}

\maketitle

En esta ayudantía, abreviaremos \emph{semilocalmente simplemente conexo} por \strong{slsc}.
Dado un grupo $(G, \cdot)$, denotaremos por $(G^{\rm op}, *)$ al \emph{grupo opuesto}, cuyos elementos son los mismos, pero cuya operación es $x*y := y\cdot x$.

% \section{Recubrimientos}
\begin{enumerate}
	\item Sea $f\colon Y \to X$ un recubrimiento de un espacio slsc $X$.
		\begin{enumerate}
			\item Pruebe que dado todo recubrimiento $Z \to Y$, la composición $Z \to X$ es un recubrimiento.
			\item\lookright
				Pruebe que, cuando $Y$ es conexo, siempre existe un único recubrimiento $Z \to Y$ tal que $Z \to X$ es un
				recubrimiento normal y, para el cual, dado todo recubrimiento normal $Z' \to Y \to X$ hay factorización
				única $Z' \to Z$.
			\item Pruebe que si $Z \to Y$ satisface que $Z \to X$ es un recubrimiento normal, entonces $Z \to Y$ también lo es.
		\end{enumerate}

	\item\lookup
		Considere el espacio $X$ dado en la fig.~\ref{fig:hawaiian_earring_a}, conocido comúnmente como los \emph{aretes hawaianos}.
		\begin{enumerate}
			\item Explique por qué el espacio $X$ no es slsc.
			\item Muestre que el espacio $Y$ (fig.~\ref{fig:hawaiian_earring_b}) es un recubrimiento de $X$, y que $Z$
				(fig.~\ref{fig:hawaiian_earring_c}) es un recubrimiento doble de $Y$.
				Sin embargo, $Z$ no es recubrimiento de $X$.
		\end{enumerate}
		\begin{figure}[!hbtp]
			\centering
			\begin{subfigure}{3cm}
				\centering
				\includegraphics{hawaiian_earring_a.pdf}
				\caption{}
				\label{fig:hawaiian_earring_a}
			\end{subfigure}
			\begin{subfigure}{7cm}
				\centering
				\includegraphics{hawaiian_earring_b.pdf}
				\caption{}
				\label{fig:hawaiian_earring_b}
			\end{subfigure}
			\begin{subfigure}{7cm}
				\centering
				\includegraphics{hawaiian_earring_c.pdf}
				\caption{}
				\label{fig:hawaiian_earring_c}
			\end{subfigure}
			\caption{}
		\end{figure}

	\item Sea $G$ un grupo topológico conexo y slsc, y sea $\pi \colon \tilde G \to G$ su recubrimiento universal.
		\begin{enumerate}
			\item Pruebe que, tras elegir $e \in \tilde G$ con $\pi(e) = 1$, existe una estructura natural de
				grupo topológico sobre $\tilde G$ para la cual $e$ es el neutro, y $\pi \colon \tilde G \to G$ es un homomorfismo continuo.
				\begin{hint}
					Puede resultar útil recordar la construcción manual del recubrimiento universal.
				\end{hint}
			\item Pruebe que la sucesión
				\[\begin{tikzcd}
					1 \rar & \pi_1(G, 1)^{\rm op} \rar & \tilde G \rar["\pi"] & G \rar & 1
				\end{tikzcd}\]
				de grupos topológicos es exacta.
		\end{enumerate}

	\item Dadas funciones continuas $f \colon Y \to X$ y $g \colon Z \to X$, se define el \emph{producto fibrado} como el subespacio
		\[
			Y \times_{f,X,g} Z = Y \times_X Z := \{ (y, z) \in Y \times Z : f(y) = g(z) \} \subseteq Y \times Z.
		\]
		\begin{enumerate}
			\item Pruebe que si $f$ es un recubrimiento, entonces $\pr_2 \colon Y\times_X Z \to Z$ también lo es.
			\item Pruebe que un recubrimiento $f\colon Y \to X$ es normal con $\Aut_X(Y) = G$ syss $\pr_1 \colon Y\times_X Y \to
				Y$ es homeomorfo (relativo a $Y$) al recubrimiento $\pr_1\colon Y \times G \to Y$, donde $G$ posee la
				topología discreta.
		\end{enumerate}

	% \item Un \strong{grupoide} es una categoría (pequeña, es decir, cuyos objetos forman un conjunto), donde todas las flechas son isomorfismos.
	% 	% Intuitivamente, un grupoide \textquote{se ve} como una unión disjunta de grupos sin una identidad distinguida.
	% 	\begin{enumerate}
	% 		\item Pruebe que darse un grupoide es lo mismo que darse un conjunto $G$ (de objetos) y un conjunto $F$ (de flechas),
	% 			con funciones $d, c \colon F \to G$ (de dominio y codominio), una función $e \colon G \to F$ (de identidad),
	% 			y una función $\iota \colon F \to F$ (de inversa), tales que
	% 	\end{enumerate}
\end{enumerate}

\section*{Comentarios adicionales}
La relación entre la teoría de Galois y la correspondencia galoisiana de recubrimientos topológicos es patente (pese a que las flechas vayan
en direcciones opuestas).
El lector podría pues considerar que hay un \textquote{yoga galoisiano} subyacente que se exhibe cuando un objeto $X$ de una categoría posee
una familia de flechas (\textquote{recubrimientos}) $Y \to X$ de modo que, por ejemplo, todo recubrimiento se extienda a uno \emph{de
Galois} $\tilde Y \to X$, donde $X$ sea el cociente de $\tilde Y$ bajo el grupo de automorfismos relativos $\Aut(Y/X)$.
Así, el $\pi_1(X)$ como grupo \textquote{sin punto base}, se recupera mediante alguna construcción basada en esta categoría de
\textquote{recubrimientos}.
Esta es una idea de Grothendieck explorada en \cite{sga-1} en el capítulo sobre categorías de Galois, y su utilidad es dar una definición
\emph{puramente algebraica} (que no depende de lazos ni métodos transcendentales) del $\pi_1$ de una variedad.\footnotemark{}
\footnotetext{O bueno, casi.
	El único problema es que, cuando permitimos recubrimientos de fibras infinitas, aparecen construcciones con funciones transcendentes no algebraicas
	(el análogo es como una función analítica puede tener infinitos ceros, como $\sin z$ por ejemplo; pero una función racional no).
	Así, el $\pi_1$ que construye Grothendieck es la \emph{compleción profinita} del $\pi_1$ topológico.
}

\printbibliography

\end{document}
